%%%%%%%%%%%%%%%%%%%%%%%%%%%%%%%%%%%%%%%%%%%%%%%%%%%%%%%%%%%%%%%%%%%%%%%%%%%%%%%%
% IMU Odometry Challenge — Technical Report Template
%%%%%%%%%%%%%%%%%%%%%%%%%%%%%%%%%%%%%%%%%%%%%%%%%%%%%%%%%%%%%%%%%%%%%%%%%%%%%%%%

\documentclass[letterpaper,10pt,conference]{ieeeconf}

\IEEEoverridecommandlockouts
\overrideIEEEmargins

% ------------------------------------------------------------------------------
% Packages
% ------------------------------------------------------------------------------
\usepackage{amsmath}
\usepackage{amssymb}
\usepackage{booktabs}
\usepackage{multirow}
\usepackage{array}
\usepackage{tabularx}
\usepackage{xcolor}
\usepackage{url}
\usepackage{hyperref}
% 链接改为着色文字而非边框；挑战赛网址按要求标红。
\hypersetup{colorlinks=true, urlcolor=red, linkcolor=black, citecolor=black}
% 打分服务的 URL 有 71 字符，双栏放不下：允许在 / - _ . 处断行。
% 不加 \Urlmuskip：那会把 URL 拉伸对齐成 "https : / / hugging face . co"，很难看。
% 也不在 '-' 处断行：断在 Tartan-IMU 的连字符上，读者会以为那个连字符是排版加的。
\def\UrlBreaks{\do\/\do\_\do\.}
\usepackage{graphicx}

% Optional colors for comments/instructions
\definecolor{instructionblue}{RGB}{35,90,160}
\newcommand{\instruction}[1]{\textcolor{instructionblue}{\textit{#1}}}

% 官方 per-sequence 打分服务（Table III/IV/V 的数字必须来自这里）。
% \urldef 在定义处冻结 catcode，因此 URL 里的下划线是安全的。
\urldef{\scoringurl}\url{https://huggingface.co/spaces/Tartan-IMU/imu_odometry_challenge_scoring}

% ------------------------------------------------------------------------------
% Title and Authors
% ------------------------------------------------------------------------------
\title{\LARGE \bf
Technical Report for the IMU Odometry Challenge
}

\author{\\
Replace with Your Submission Title\\
Kaggle team name: \textcolor{red}{\textit{exactly as it appears on the leaderboard}}\\
Albert Author$^{1}$, Bernard D. Researcher$^{2}$, and Charlie Contributor$^{3}$%
\thanks{Replace this footnote with funding or support acknowledgments, if applicable.}%
\thanks{$^{1}$Albert Author is with Department/Institution, City, Country.
        {\tt\small albert.author@example.com}}%
\thanks{$^{2}$Bernard D. Researcher is with Department/Institution, City, Country.
        {\tt\small bernard.researcher@example.com}}%
\thanks{$^{3}$Charlie Contributor is with Department/Institution, City, Country.
        {\tt\small charlie.contributor@example.com}}%
}

\begin{document}

\maketitle
\thispagestyle{empty}
\pagestyle{empty}

%%%%%%%%%%%%%%%%%%%%%%%%%%%%%%%%%%%%%%%%%%%%%%%%%%%%%%%%%%%%%%%%%%%%%%%%%%%%%%%%
\begin{abstract}
\instruction{Please read the latest notice from
\url{https://superodometry.com/imuchallenge}
\\
Provide a concise summary of your solution to the IMU Odometry Challenge.
The abstract should briefly describe:
(1) the core idea and overall approach;
(2) the main model or system architecture;
(3) the most important techniques used to improve performance;
(4) the training data and inference strategy, when relevant; and
(5) the final challenge results using the primary evaluation metrics.
}

\end{abstract}

\vspace{0.5em}
\noindent\textcolor{red}{\rule{\columnwidth}{0.4pt}}\\[-0.3em]
\noindent\textcolor{red}{\textbf{Selection criteria.} The Kaggle leaderboard is \emph{not} the
only criterion used by the challenge committee to determine the Top 10 teams: the quality of
this technical report is also taken into account. The Top 10 teams will be invited to
contribute to the forthcoming \emph{IMU Foundation Model} white paper.}\\[-0.2em]
\noindent\textcolor{red}{\textbf{Format and deadline.} Maximum \textbf{6 pages excluding
references}, and \textbf{7 pages in total} including references, in this template. The
\textbf{appendix does not count} towards either limit. Reports are due
\textbf{23:59 US Eastern Time (EDT, UTC$-$4) on 23 September 2026}, three days after the
leaderboard closes.}\\[-0.2em]
\noindent\textcolor{red}{\textbf{How to submit.} The report is \emph{not} uploaded on the
website. It is submitted \textbf{inside the challenge Form}: fill the Form in as usual and
attach this report at the end --- that single submission is what we collect.
The Form is linked from the \textbf{Announcements} section of
\url{https://superodometry.com/imuchallenge}. Please check that section regularly: all
challenge news, rule clarifications and any change to these deadlines are posted there.}\\[-0.2em]
\noindent\textcolor{red}{\textbf{Where the numbers in this report come from.} Every number in
Tables~\ref{tab:overall_results}, \ref{tab:platform_results} and \ref{tab:sequence_results}
must be produced by the \textbf{official per-sequence scoring service}:\\[0.15em]
\scoringurl\\[0.15em]
Upload your submission CSV and it returns ATE$_{20}$, AVE and RTE for all 89 test sequences,
computed with the same code as the Kaggle leaderboard. \textbf{Five submissions per team per
day.} Your team name must match your \textbf{Kaggle team name exactly}; teams with no Kaggle
submission are not scored.}\\[-0.2em]
\noindent\textcolor{red}{\rule{\columnwidth}{0.4pt}}
\vspace{0.5em}

%%%%%%%%%%%%%%%%%%%%%%%%%%%%%%%%%%%%%%%%%%%%%%%%%%%%%%%%%%%%%%%%%%%%%%%%%%%%%%%%
\section{Introduction}

\instruction{
Briefly introduce your approach to the challenge.
We recommend addressing the following questions:
}

\begin{itemize}
    \item What are the main challenges of developing foundation models for IMU odometry?
    \item What motivated the design of your solution?
    \item What are the key components of your method?
    \item What distinguishes your method from others?
    \item What are the main contributions of your method?
    
    %\item What are the main contributions of your challenge submission?
    
    % \item What assumptions does your method make about the IMU, platform,
    % motion distribution, coordinate frame, or deployment setting?
\end{itemize}

%%%%%%%%%%%%%%%%%%%%%%%%%%%%%%%%%%%%%%%%%%%%%%%%%%%%%%%%%%%%%%%%%%%%%%%%%%%%%%%%
\section{Method}

\instruction{
Describe the proposed method in sufficient detail to allow readers to understand,
analyze, and reproduce the system.
}

% ------------------------------------------------------------------------------
\subsection{Overall Architecture}

\begin{itemize}
    \item What is the overall system architecture of the proposed solution? 
    % \item What are the model inputs and outputs?
    % \item How is temporal IMU information represented and modeled?
    % \item What are the main network components?
    % \item Is a single unified model used for all platforms, or are
    % platform-specific models/components used?
    %\item Does the system use routing, mixture-of-experts, ensembling, hierarchical prediction, recurrent state, or multi-stage refinement?
    %\item Does the system explicitly estimate intermediate quantities such as velocity, displacement, orientation, bias, gravity, uncertainty, or motion class?
    \item What design choices were made specifically to improve cross-platform
    or out-of-distribution generalization?
\end{itemize}

\instruction{
Please include an architecture or pipeline figure.
}

\begin{figure}[htbp]
    \centering
    % \includegraphics[width=\linewidth]{figures/method_overview.pdf}
    \fbox{\parbox[c][3.2cm][c]{0.92\linewidth}{
        \centering
        Replace with method overview / architecture figure
    }}
    \caption{Overview of the proposed IMU odometry method.}
    \label{fig:method_overview}
\end{figure}

% ------------------------------------------------------------------------------
\subsection{Data Processing}

\begin{itemize}
    %\item What raw IMU signals are used?
    %\item What sampling frequency is used during training and inference?
    %\item Is resampling or synchronization required?
    %\item How are accelerometer and gyroscope measurements normalized?
    %\item Are coordinate-frame transformations applied?
    %\item Is gravity alignment, orientation normalization, or canonicalization used?
    %\item How are training windows or temporal sequences constructed?
    %\item What temporal context length and stride are used?
    %\item Are sensor biases, offsets, or calibration parameters estimated or corrected?
    \item What is the data preprocessing process?
        
    %\item What filtering, clipping, denoising, interpolation, or missing-data handling is used?
    \item What data augmentation (if any) is applied?
    %\item Are augmentations physically motivated, e.g., rotation, bias, scale, noise, time warping, sensor perturbation, or trajectory perturbation?
    \item Anything else? 
\end{itemize}

% ------------------------------------------------------------------------------
\subsection{Training Data}

\instruction{
Clearly list \textbf{all datasets} used to develop or train the final submitted model.
This includes official challenge data, external public datasets, synthetic data,
self-collected data, pretrained checkpoints, and pseudo-labeled data.
}

\begin{table}[htbp]
    \centering
    \footnotesize
    \caption{Datasets used for model development and training.}
    \begin{tabular}{p{2.3cm} p{1.6cm} p{2.8cm}}
        \toprule
        Dataset & Type & Usage \\
        \midrule
        Challenge train set & Official & Training / validation \\
        External Dataset A & Public & Pretraining \\
        Synthetic Dataset B & Synthetic & Augmentation \\
        \bottomrule
    \end{tabular}
    
    \label{tab:datasets}
    
\end{table}

% \begin{itemize}
%     \item Which datasets were used for training, validation, model selection,
%     pretraining, or hyperparameter tuning?
%     \item Were any external datasets used? If so, why were they useful?
%     \item Was synthetic data used? If so, how was it generated?
%     \item Were pretrained models or released checkpoints used?
%     \item How were train/validation splits constructed?
%     \item How did you avoid test-set leakage or unintended use of test information?
%     \item Did you balance or reweight platforms, sequences, or motion types during training?
% \end{itemize}

% ------------------------------------------------------------------------------
\subsection{Training and Optimization}

\begin{itemize}
    \item What loss functions, optimizer, learning-rate schedule, batch size  are used?
    % \item What are the major hyperparameters?
    \item How many epochs or optimization steps are used?  What hardware is used for training? And what is the approximate training time?
    
    \item Is training performed in one stage or multiple stages?
    \item Is curriculum learning, fine-tuning, domain adaptation, or platform-specific specialization used?
    \item How is the final checkpoint selected?
    % \item Are multiple random seeds used?
    % \item Are exponential moving average, checkpoint averaging, or other
    % stabilization techniques used?
    % \item What regularization techniques are used?
\end{itemize}





\subsection{Inference}

\begin{itemize}
    \item Describe the complete inference pipeline used to generate the final submission.

    \item What is the average inference latency or throughput?
    % \item Does inference use checkpoint ensembling or model ensembling?
    % \item Is test-time augmentation (TTA) used?
    % \item Is temporal smoothing or trajectory-level post-processing used?
    % \item Is platform classification, routing, or specialization used?
    % \item Does inference maintain persistent hidden state?
    % \item Does the method use future context, full-trajectory context, or only past observations?
    % \item Is any test-time / online adaptation performed?
    % \item Are predictions integrated, filtered, calibrated, or otherwise post-processed?
    % \item If the development-time system differs from the submitted system,
    % what exactly changed in the final submission?
\end{itemize}


% ------------------------------------------------------------------------------
\subsection{Model Complexity and Runtime}

\begin{itemize}
    \item How many parameters does the final model contain?
   
    \item What hardware is used for inference?

    \item What is your inference speed? (Hz)
  
    \item Can the method operate online / in real time?
    \item What is the peak GPU/CPU memory requirement?
    % \item Does the method require the full trajectory or only causal/past context?
\end{itemize}

\begin{table}[htbp]
    \centering
    \footnotesize
     \caption{Model complexity and runtime statistics.}
    \begin{tabular}{lc}
        \toprule
        Property & Value \\
        \midrule
        Number of parameters & -- \\
        Training hardware & -- \\
        Training time & -- \\
        Inference hardware & -- \\
        Inference speed (Hz) & -- \\
        % Temporal context & -- \\
        Peak memory & -- \\
        \bottomrule
    \end{tabular}
   
    \label{tab:runtime}
\end{table}

% ------------------------------------------------------------------------------

%%%%%%%%%%%%%%%%%%%%%%%%%%%%%%%%%%%%%%%%%%%%%%%%%%%%%%%%%%%%%%%%%%%%%%%%%%%%%%%%
\section{Development Experiments}
\label{sec:devexp}

\instruction{
Summarize \textbf{what you actually tried and what happened} --- the experimental record.
Interpretation and lessons for the field belong in Section~\ref{sec:insights}, not here.
Negative results are encouraged when they provide useful technical insight; please quantify
gains rather than describing them qualitatively.
}

\begin{itemize}
    \item What model architectures did you try?
    \item Which changes produced the largest gains, and \textbf{by how much} on which metric?
    \item Which approaches did not work as expected?
    % \item What experiments changed your understanding of the problem?
    % \item Did improvements transfer consistently across platforms?
    \item Were some techniques beneficial for one platform but harmful for another?
    \item Which components improved generalization rather than only validation performance?
    % \item Did additional temporal context consistently help?
    \item Did more training data consistently help?
    \item Did model scale correlate with better performance?
    % \item What was the most important bottleneck: model capacity, data quality, representation, optimization, domain shift, or inference strategy?
\end{itemize}

%%%%%%%%%%%%%%%%%%%%%%%%%%%%%%%%%%%%%%%%%%%%%%%%%%%%%%%%%%%%%%%%%%%%%%%%%%%%%%%%
\section{Experiments}

% ------------------------------------------------------------------------------
\subsection{Official Challenge Performance}

\instruction{
Report the performance of the final submitted method using the official challenge
evaluation metrics, and include the leaderboard result or ranking.
}

\noindent\textcolor{red}{\textbf{All numbers in Tables~\ref{tab:overall_results},
\ref{tab:platform_results} and \ref{tab:sequence_results} must be produced by the official
scoring service, \scoringurl, not computed locally.} State the MD5 of
the submission CSV they come from, so that the committee can reproduce them:
\texttt{md5 = \textbf{[fill in]}}.}

\begin{table}[htbp]
    \centering
    \footnotesize
    \caption{Overall challenge performance on the \textbf{final (private) leaderboard}.
    The TartanIMU Score is $0.6\,\mathrm{AVE}/0.7356 + 0.4\,\mathrm{ATE}_{20}/3.1160$;
    it is dimensionless and \emph{lower is better} (the all-zero submission scores exactly
    1.000). ATE$_{20}$ is the 20\,m-segment trajectory error after SE(3) alignment.
    \textbf{RTE does not enter the ranking} and is reported for analysis only; it is defined
    in the caption of Table~\ref{tab:sequence_results}.}
    \setlength{\tabcolsep}{3.5pt}
    \begin{tabular}{lcccc}
        \toprule
        Method & Score $\downarrow$ & ATE$_{20}$ $\downarrow$ & AVE $\downarrow$ & RTE $\downarrow$ \\
        \midrule
        TartanIMU Baseline & -- & -- & -- & -- \\
        Ours & -- & -- & -- & -- \\
        \bottomrule
    \end{tabular}
    
    \label{tab:overall_results}
\end{table}

% ------------------------------------------------------------------------------
\subsection{Per-Platform Results}

\begin{table}[htbp]
    \centering
    \footnotesize
    \caption{Performance by platform. RTE is defined in the caption of
    Table~\ref{tab:sequence_results}.}
    \begin{tabular}{lccc}
        \toprule
        Platform & ATE$_{20}$ $\downarrow$ & AVE $\downarrow$ & RTE $\downarrow$ \\
        \midrule
        Wheeled / Car & -- & -- & -- \\
        Handheld / Human & -- & -- & -- \\
        Legged / Quadruped & -- & -- & -- \\
        Aerial / Drone & -- & -- & -- \\
        \midrule
        Average & -- & -- & -- \\
        \bottomrule
    \end{tabular}
    
    \label{tab:platform_results}
\end{table}

% ------------------------------------------------------------------------------
\subsection{Per-Sequence Results}

\instruction{The full per-sequence table is reported in the appendix,
Table~\ref{tab:sequence_results}. Please discuss the trends here rather than repeating the numbers:
which sequences are hardest, what they have in common, and whether the errors are concentrated
in a few sequences or spread evenly.}

% ------------------------------------------------------------------------------
% \subsection{Comparison with the TartanIMU Baseline}

% \begin{itemize}
%     \item How much does the final method improve over the provided TartanIMU baseline?
%     \item Which platforms show the largest gains?
%     \item Which platforms remain the most difficult?
%     \item Which components account for the largest improvements?
%     % \item Are improvements consistent across sequences, or driven by a small subset?
% \end{itemize}

% ------------------------------------------------------------------------------
\subsection{Ablation Studies}

\instruction{
Ablation studies are strongly encouraged for the most important components
of the final method.
}



% Suggested ablations include:
% \begin{itemize}
%     \item model architecture;
%     \item temporal context length;
%     \item input representation;
%     \item data augmentation;
%     \item external or synthetic training data;
%     \item loss functions;
%     \item model scaling;
%     \item platform specialization or routing;
%     \item ensembling;
%     \item test-time augmentation;
%     \item trajectory post-processing.
% \end{itemize}

% ------------------------------------------------------------------------------
\subsection{Generalization Analysis}

\begin{itemize}
    \item How well does the method generalize?
    \item Is there evidence of negative transfer between platforms?
    % \item Does one unified model perform competitively with platform-specific models?
    % \item Which motion types or platforms benefit most from shared training?
    % \item What characteristics of the training distribution appear most important
    % for test-time generalization?
    % \item How sensitive is the method to sampling rate, sensor noise, bias,
    % orientation, or coordinate-frame differences?
\end{itemize}

% ------------------------------------------------------------------------------
\subsection{Failure-Mode Analysis}

\begin{itemize}
    \item What are the most common failure cases?
    % \item Does performance degrade under aggressive motion or long-duration integration?
    % \item How does the method behave under unseen motion patterns?
    % \item How sensitive is it to sensor bias, noise, saturation, or calibration errors?
    % \item Are failures associated with particular platforms or motion regimes?
    % \item Can failures be detected from model confidence or uncertainty?
    \item What limitations remain unresolved?
\end{itemize}

%%%%%%%%%%%%%%%%%%%%%%%%%%%%%%%%%%%%%%%%%%%%%%%%%%%%%%%%%%%%%%%%%%%%%%%%%%%%%%%%
\section{Insights and Lessons Learned}
\label{sec:insights}

\instruction{
This section is about \textbf{what your results mean for the field}, not about what you ran
--- the experimental record belongs in Section~\ref{sec:devexp}. Please interpret those
numbers rather than repeating them.
}

\begin{itemize}
    \item Based on your experiments, does a \textbf{single universal inertial model} appear
    feasible, or does platform specialization remain necessary? What evidence points either way?
    \item Where does the bottleneck currently sit --- data, architecture, optimization, or the
    evaluation protocol itself?
    \item Which of your findings would you expect to \textbf{transfer} to platforms, sensors or
    motion regimes outside this benchmark, and which would not?
    \item What kinds of data gave the highest marginal benefit per hour collected?
    \item What did this challenge reveal about the current limitations of learning-based
    inertial odometry?
    \item Which research directions look most promising for future IMU odometry systems?
    % \item What types of data provided the highest marginal benefit?
    % \item Is a universal inertial model feasible based on your experiments?
    % \item When does platform-specific specialization appear necessary?
    % \item What did the challenge reveal about the current limitations of learning-based inertial odometry?
    % \item Based on your experiments, what research directions appear most promising for future IMU odometry systems?
\end{itemize}

%%%%%%%%%%%%%%%%%%%%%%%%%%%%%%%%%%%%%%%%%%%%%%%%%%%%%%%%%%%%%%%%%%%%%%%%%%%%%%%%

% ------------------------------------------------------------------------------
\subsection{Feedback to the Organizers}

\instruction{This subsection is read directly by the challenge committee and feeds into the
design of the next edition. Candid criticism is more useful to us than politeness.}

\begin{itemize}
    \item Were the rules, metrics, datasets and submission process clear and unambiguous?
    Which part cost you the most time to figure out?
    \item Did anything about the evaluation metric behave in a way you did not expect?
    \item What additional data, platforms, or evaluation settings would most improve the
    next edition?
    \item Was any part of the benchmark measuring something other than what it intended to?
\end{itemize}

%%%%%%%%%%%%%%%%%%%%%%%%%%%%%%%%%%%%%%%%%%%%%%%%%%%%%%%%%%%%%%%%%%%%%%%%%%%%%%%%
\section{Conclusion}
\instruction{Summarize the overall report: the final solution and its main strengths, the
most important remaining limitations, and the most promising future directions.}

\section*{ACKNOWLEDGMENT}
\instruction{Please include the acknowledgment here.}

%%%%%%%%%%%%%%%%%%%%%%%%%%%%%%%%%%%%%%%%%%%%%%%%%%%%%%%%%%%%%%%%%%%%%%%%%%%%%%%%
\section*{APPENDIX}
\instruction{Add the appendix here.}

% The appendix may include:
% \begin{itemize}
%     \item additional implementation details;
%     \item full hyperparameter settings;
%     \item extended ablation studies;
%     \item additional per-sequence tables;
%     \item qualitative examples;
%     \item additional failure cases;
%     \item training curves;
%     \item supplementary reproducibility information.
% \end{itemize}


\subsection{Per-Sequence Results (Full Table)}

\instruction{Report the metrics for \textbf{all 89 test sequences}. Upload your submission
CSV to the official scoring service, \scoringurl, and it returns ATE, AVE and RTE for
every sequence, ready to be copied into the table below (the service also offers the same
table as a CSV download).}

\begin{table*}[htbp]
    \centering
    \footnotesize
    \caption{Per-sequence challenge results over the 89 test sequences. ATE, AVE and RTE are
    produced by the official per-sequence scoring service whose URL is given on the first page.
    \textbf{RTE} (Relative Trajectory Error, m) follows the definition of Sturm
    \emph{et al.}~\cite{sturm2012rgbd} and is the metric used in AirIO: over sliding windows of
    a fixed duration $\Delta t$, it is the RMSE of the difference between the predicted and the
    ground-truth relative displacement,
    $\mathrm{RTE} = \sqrt{\frac{1}{n}\sum_{i}\lVert (\hat{\mathbf{p}}_{i+\Delta t} -
    \hat{\mathbf{p}}_{i}) - (\mathbf{p}_{i+\Delta t} - \mathbf{p}_{i}) \rVert^{2}}$.
    Unlike ATE it measures \emph{local} drift and is insensitive to accumulated global offset.
    The window length is fixed to $\Delta t = 5$\,s for all platforms: at $\Delta t = 1$\,s the
    metric is nearly a restatement of AVE, while the shortest aerial sequences are only 19\,s
    long, which bounds $\Delta t$ from above.}
    \label{tab:sequence_results}
    \begin{tabular}{cccc @{\hspace{2.2em}} cccc @{\hspace{2.2em}} cccc}
        \toprule
        Seq. ID & ATE$_{20}$ $\downarrow$ & AVE $\downarrow$ & RTE $\downarrow$ &
        Seq. ID & ATE$_{20}$ $\downarrow$ & AVE $\downarrow$ & RTE $\downarrow$ &
        Seq. ID & ATE$_{20}$ $\downarrow$ & AVE $\downarrow$ & RTE $\downarrow$ \\
        \midrule
        1  & -- & -- & -- & 31 & -- & -- & -- & 61 & -- & -- & -- \\
        2  & -- & -- & -- & 32 & -- & -- & -- & 62 & -- & -- & -- \\
        3  & -- & -- & -- & 33 & -- & -- & -- & 63 & -- & -- & -- \\
        $\vdots$ & & & & $\vdots$ & & & & $\vdots$ & & & \\
        30 & -- & -- & -- & 60 & -- & -- & -- & 89 & -- & -- & -- \\
        \bottomrule
    \end{tabular}
\end{table*}

% ------------------------------------------------------------------------------
\subsection{Reproducibility and Submission Artifacts}

\instruction{
List every artifact needed to reproduce the reported numbers: Kaggle team name and the
submission that produced them (submission id and MD5 of the CSV); code repository and commit;
trained checkpoint (link + MD5); configuration files; environment;
random seeds, and whether the reported
score is a single run or an average.
}

\vspace{0.3em}
\noindent\textcolor{red}{\textbf{Compliance declaration (required).}} \emph{Every team must
answer all five questions below explicitly. Reports that leave them unanswered cannot be
considered for the Top 10.}

\begin{enumerate}
    \item Do the submitted predictions come from a \textbf{single model with one shared set of
    weights}? Inferring the embodiment inside the network is allowed and is the point of the
    benchmark; dispatching to separately trained per-platform experts is not.
    \item Is any \textbf{platform classification, routing, or platform-specific
    specialization} used at inference? If yes, describe exactly where it sits in the pipeline.
    \item Is any \textbf{ensembling} used --- model ensembling, checkpoint averaging or
    soups, or test-time augmentation? If yes, state how many members and how they are combined.
    \item How did you \textbf{avoid test-set leakage}, including any indirect use of test
    information for model or checkpoint selection?
    \item If the development-time system differs from the submitted system, \textbf{what
    exactly changed} in the final submission?
\end{enumerate}

% Please specify:
% \begin{itemize}
%     \item Kaggle team name, and the exact submission that produced the
%           reported score (submission id / filename / MD5);
%     \item code repository and commit hash;
%     \item trained checkpoint (download link + MD5);
%     \item configuration file(s) used for the final inference run;
%     \item environment specification (Python / PyTorch / CUDA versions);
%     \item hardware and wall-clock time, for training and for inference;
%     \item random seeds, and whether the reported score is a single run or an
%           average over seeds --- if a single run, state the run-to-run spread;
%     \item \textbf{rule compliance}: confirm the predictions come from a
%           \emph{single model with one shared set of weights}. Inferring the
%           embodiment inside the network is allowed and is the point of the
%           benchmark; recovering the platform in order to dispatch to four
%           separately trained experts is not;
%     \item any external data or pretrained weights used beyond the released
%           TartanIMU baseline.
% \end{itemize}

% ------------------------------------------------------------------------------
\subsection{Final Predictions}

\instruction{
Provide a link to the final predictions on the challenge test set. This must be
the same CSV that was submitted to Kaggle: exactly the columns below, one row
per test \texttt{window\_id}, 30,644 rows plus the header.
}

\begin{verbatim}
window_id,vx,vy,vz
1,0.0,0.0,0.0
2,0.0,0.0,0.0
...
\end{verbatim}

\noindent
\texttt{vx, vy, vz} is the predicted mean \textbf{body-frame} velocity, in m/s,
of that window. A window is \texttt{1.0 s = 200 frames} at 200\,Hz,
non-overlapping: window $k$ of a trajectory is \texttt{imu[k*200:(k+1)*200]}.
Window ids are taken from \texttt{index/test\_windows.csv}; every id must appear
exactly once.

% Please specify:
% \begin{itemize}
%     \item prediction download link;
%     \item which Kaggle submission this file corresponds to;
%     \item MD5 of the submitted CSV;
%     \item any post-processing applied after the network output
%           (test-time augmentation, temporal smoothing, magnitude rescaling,
%           ensembling or model soup);
%     \item how degenerate or missing windows are handled, if any.
% \end{itemize}

% Example:
% \url{https://drive.google.com/...}

% ------------------------------------------------------------------------------
\subsection{Model and Inference Code}

\instruction{Please provide final trained model/checkpoint and inference script.}

% Please provide, when possible:
% \begin{itemize}
%     \item final trained model/checkpoint;
%     \item inference script;
%     \item configuration files required for inference;
%     \item environment/dependency specification;
%     \item code repository;
%     \item brief instructions for reproducing the submitted predictions.
% \end{itemize}

% \instruction{
% If the code or model cannot be publicly released, please describe the reason
% and provide enough implementation detail for readers to understand the system.
% }


%%%%%%%%%%%%%%%%%%%%%%%%%%%%%%%%%%%%%%%%%%%%%%%%%%%%%%%%%%%%%%%%%%%%%%%%%%%%%%%%


%%%%%%%%%%%%%%%%%%%%%%%%%%%%%%%%%%%%%%%%%%%%%%%%%%%%%%%%%%%%%%%%%%%%%%%%%%%%%%%%
% References
%%%%%%%%%%%%%%%%%%%%%%%%%%%%%%%%%%%%%%%%%%%%%%%%%%%%%%%%%%%%%%%%%%%%%%%%%%%%%%%%

\begin{thebibliography}{99}

\bibitem{tartanimu}
S.~Zhao, S.~Zhou, R.~Blanchard, Y.~Qiu, W.~Wang, and S.~Scherer,
``Tartan IMU: A Light Foundation Model for Inertial Positioning in Robotics,''
in \emph{Proc.\ IEEE/CVF Conf.\ on Computer Vision and Pattern Recognition
(CVPR)}, 2025, pp.~22520--22529.

\bibitem{sturm2012rgbd}
J.~Sturm, N.~Engelhard, F.~Endres, W.~Burgard, and D.~Cremers,
``A Benchmark for the Evaluation of RGB-D SLAM Systems,''
in \emph{Proc.\ IEEE/RSJ Int.\ Conf.\ on Intelligent Robots and Systems (IROS)}, 2012,
pp.~573--580.

\bibitem{baseline2}
Replace with additional references as needed.

\end{thebibliography}

\end{document}